\chapter{Scénarios de test}\label{chap:tests}

\section{Stratégie}
Trois niveaux de tests complémentaires ont été mis en place :
\begin{enumerate}
  \item \textbf{Tests d'intégration automatisés} (JUnit 5, \code{PlateformeIT}) : la plateforme complète est démarrée
        sur des bases H2 en mémoire puis sollicitée par de \emph{vrais} appels SOAP (clients JAX-WS générés depuis le WSDL)
        et HTTP. Les tests couvrent les services un par un puis le processus de bout en bout, y compris les pannes.
        Ils s'exécutent à chaque \code{mvn package} et ont également été exécutés sur \textbf{PostgreSQL~16}.
  \item \textbf{Collection Postman} (\code{docs/postman}) : 39 requêtes enchaînées avec 45 assertions sur les codes
        et les contenus, authentifiées par clé d'API, exécutable dans le \emph{Collection Runner} ou avec Newman.
  \item \textbf{Requêtes SoapUI} (\code{docs/soap}) : enveloppes prêtes à l'emploi pour chaque type d'opération,
        accompagnées des réponses réellement obtenues.
\end{enumerate}

\section{Tests automatisés}
\begin{longtable}{@{}lp{7.2cm}p{5.2cm}c@{}}
\caption{Scénarios de test automatisés et résultats}\label{tab:tests}\\
\toprule
\textbf{ID} & \textbf{Scénario} & \textbf{Résultat attendu} & \textbf{Résultat} \\
\midrule
\endfirsthead
\toprule
\textbf{ID} & \textbf{Scénario} & \textbf{Résultat attendu} & \textbf{Résultat} \\
\midrule
\endhead
\bottomrule
\endfoot
\multicolumn{4}{@{}l}{\textit{Services SOAP}}\\
ST-01 & Récupérer les trois WSDL & Opérations présentes dans chaque contrat & \ok \\
ST-02 & Appel SOAP sans en-tête WS-Security & Fault \code{FailedAuthentication} & \ok \\
ST-03 & Ajouter, consulter, modifier, supprimer un étudiant & Statut \code{ACTIF} par défaut ; consultation après suppression : \code{RESSOURCE\_INTROUVABLE} & \ok \\
ST-04 & Ajouter un étudiant (matricule, e-mail et niveau invalides) & Fault \code{DONNEES\_INVALIDES} avec 3 détails & \ok \\
ST-05 & Ajouter un matricule existant & Fault \code{CONFLIT} & \ok \\
ST-06 & Un étudiant consulte une autre fiche / supprime sa fiche & Fault \code{ACCES\_REFUSE} & \ok \\
ST-07 & Vérifier un étudiant actif, suspendu, inexistant & éligible / non éligible / inexistant & \ok \\
ST-08 & CRUD entreprise ; offres d'une entreprise ; suppression avec offres actives & 2 offres ; \code{REGLE\_METIER} & \ok \\
ST-09 & Publier une offre : entreprise non référencée, autre entreprise, entreprise propriétaire & \code{REGLE\_METIER}, \code{ACCES\_REFUSE}, offre \code{ACTIVE} & \ok \\
ST-10 & Compatibilité AMOA/INE3, DATA/INE1 et offre expirée & compatible, score 50 \% ; incompatible, 2 motifs ; incompatible (date limite) & \ok \\
ST-11 & Disponibilité d'une offre expirée & date dépassée, non disponible & \ok \\
ST-12 & Fiche d'un candidat consultée par une entreprise puis par le responsable & téléphone et e-mail masqués, adresse absente ; fiche complète & \ok \\
\multicolumn{4}{@{}l}{\textit{Services REST}}\\
RT-01 & Requêtes sans identifiants, avec mauvais mot de passe, rôle étudiant sur \code{POST /stages} & 401, 401, 403 & \ok \\
RT-02 & Dépôt invalide, valide, en double, JSON mal formé & 400, 201 \code{DEPOSEE}, 409, 400 & \ok \\
RT-03 & Transition interdite, annulation par un tiers, passage en évaluation, annulation refusée, historique ; annulation d'une candidature déposée & 409, 403, 200, 422, 3 entrées ; \code{ANNULEE} puis 422 & \ok \\
RT-04 & Notification d'acceptation ; lecture par un autre étudiant ; type inconnu & 201 et message rendu ; 403 ; 400 & \ok \\
RT-05 & Stage : création, doublon, compte rendu, note 25, évaluation finale, historique & 201 \code{EN\_COURS}, 409, 201, 400, \code{TERMINE}, moyenne 16,5, étudiant \code{ACTIF} & \ok \\
RT-06 & Clé d'API valide, invalide, régénérée ; ancienne clé & 200 ; 401 \code{CLE\_API\_INVALIDE} ; 200 ; 401 (révoquée) & \ok \\
\multicolumn{4}{@{}l}{\textit{Processus BPMN}}\\
PT-01 & Scénario nominal : candidature puis acceptation par l'entreprise & Stage créé aux dates saisies, étudiant \code{EN\_STAGE}, offre \code{POURVUE}, 3 notifications ; tâche invisible pour une autre entreprise & \ok \\
PT-02 & Candidatures non recevables : étudiant suspendu, niveau insuffisant, date dépassée, entreprise non référencée, offre pourvue & 5 instances terminées par « Candidature rejetée » avec motif & \ok \\
PT-03 & Complément sans commentaire, complément, puis refus & 400 ; retour à l'évaluation ; \code{REFUSEE} ; notifications \code{DEMANDE\_COMPLEMENT} et \code{REFUS} & \ok \\
PT-04 & Service Candidatures en panne pendant les vérifications, puis relance & \code{INCIDENT\_TECHNIQUE} (étape \code{VERIFICATIONS}) puis \code{EN\_ATTENTE\_DECISION} & \ok \\
PT-05 & Service Stages en panne pendant la finalisation, puis abandon & \code{INCIDENT\_TECHNIQUE} (étape \code{FINALISATION}) puis \code{ABANDONNEE} & \ok \\
PT-06 & Service Notifications en panne pendant le dépôt & Processus poursuivi, incident consigné & \ok \\
PT-07 & Annulation d'une candidature déposée & Instance \code{INTERROMPU} & \ok \\
PT-08 & Délai de décision dépassé (2 s en test) & Relance de l'entreprise, évaluation toujours en attente & \ok \\
PT-09 & Prise en charge par l'entreprise puis tentative d'annulation & \code{DEPOSEE} → \code{EN\_COURS\_EVALUATION}, tâche attribuée, annulation refusée (422) & \ok \\
PT-10 & Sixième candidature en cours pour un même étudiant & Rejet motivé par le processus ; 422 sur l'API & \ok \\
\end{longtable}

\begin{lstlisting}[language=sh, caption={Exécution de la suite (extrait de la sortie Maven Surefire)}]
$ mvn test
Tests run: 12, Failures: 0, Errors: 0, Skipped: 0 -- in PlateformeIT$Soap
Tests run: 6,  Failures: 0, Errors: 0, Skipped: 0 -- in PlateformeIT$Rest
Tests run: 10, Failures: 0, Errors: 0, Skipped: 0 -- in PlateformeIT$Processus

$ mvn test -Ddb.url=jdbc:postgresql://localhost:5433/stages -Dbpmn.db.url=jdbc:postgresql://...
Tests run: 24, Failures: 0, Errors: 0, Skipped: 0
\end{lstlisting}

\begin{lstlisting}[language=Java, caption={Exemple de test : indisponibilité puis relance (PT-04)}]
@Test
@DisplayName("PT-04 Indisponibilité d'un service : incident, tâche du responsable, relance réussie")
void incidentRelance() {
    SimulateurPannes.definir("candidatures", true);
    JsonNode r;
    try {
        r = soumettre("E2026003", 4).json();
    } finally {
        SimulateurPannes.definir("candidatures", false);
    }
    assertEquals("INCIDENT_TECHNIQUE", r.get("statut").asText());
    assertEquals("VERIFICATIONS", r.get("variables").get("etapeCourante").asText());
    JsonNode apres = Http.post("/api/processus/taches/" + tache(r, "traiterIncident") + "/incident",
            ADMIN, ADMIN_PWD, Map.of("action", "RELANCER")).json();
    assertEquals("EN_ATTENTE_DECISION", apres.get("statut").asText());
}
\end{lstlisting}

\section{Collection Postman}
La collection est organisée en six dossiers (système, processus, candidatures, notifications, stages, pannes). Les
variables de collection (\code{tacheId}, \code{candidatureId}, \code{stageId}, \code{incidentId}) sont alimentées par
les scripts de test, ce qui permet d'enchaîner un scénario complet.

\begin{lstlisting}[language=sh, caption={Exécution de la collection avec Newman}]
$ newman run docs/postman/gestion-stages.postman_collection.json
|                requests |                 39 |                 0 |
|              assertions |                 45 |                 0 |
\end{lstlisting}

\section{Requêtes SoapUI}
Après import des trois WSDL dans SoapUI, les enveloppes de \code{docs/soap/} peuvent être collées telles quelles :

\begin{table}[H]
\centering\small
\caption{Requêtes SOAP fournies}
\begin{tabularx}{\textwidth}{@{}lYl@{}}
\toprule
\textbf{Fichier} & \textbf{Cas} & \textbf{Résultat obtenu} \\
\midrule
\code{01-etudiants-ajouterEtudiant} & Création d'une étudiante & fiche \code{ACTIF} \\
\code{02-etudiants-consulterEtudiant} & Consultation & fiche complète \\
\code{03-etudiants-verifierEtudiant} & Étudiante suspendue & \code{eligible=false} \\
\code{04-etudiants-ajouterEtudiant-invalide} & Données invalides & Fault, 3 détails \\
\code{05-entreprises-verifierEntreprise} & Entreprise non référencée & \code{referencee=false} \\
\code{06-entreprises-consulterOffresEntreprise} & Composition SOAP & 2 offres \\
\code{07-offres-creerOffre} & Publication d'une offre & offre \code{ACTIVE} \\
\code{08-offres-verifierDisponibilite} & Offre expirée & \code{disponible=false} \\
\code{09-offres-verifierCompatibilite} & Score de compétences & 50 \%, 2 manques \\
\code{10-sans-ws-security} & Sans authentification & \code{wsse:FailedAuthentication} \\
\bottomrule
\end{tabularx}
\end{table}
